% Insercion editor-ready para el capitulo 31.

\section{Dos lectores internos de \texorpdfstring{\(\alpha\)}{alpha}: clausura dodecafásica y núcleo analítico nonádico de orden \texorpdfstring{\(9\)}{9}}
\label{sec:l9s102:alpha-dos-vias}

Esta sección desarrolla la coordenada \(\alpha\) de la
\hyperref[sec:tesis-inaugural-helicoidal-hmt-md]{tesis inaugural}. Ambos
lectores pertenecen a la misma genealogía APP--TRIT--TPK y preceden a la
comparación metrológica. La vía A opera mediante la clausura dodecafásica
reversible del estado firmado y publica la sección coinductiva
\(\alpha_{\mathrm{HMT}}\). La vía B construye, sin consultar esa salida, el
núcleo analítico nonádico exacto de orden nueve
\(E_{21/22}^{(9)}:=P_9-\log_{10}(10/9)\) y publica su raíz positiva simple y
única \(\xi_9\). Sus dominios, operaciones y certificados se mantienen
diferenciados:
\begin{equation}
 X_{12}^{\mathrm{sig}}
 \xrightarrow{\ \Gamma_{12}^{\alpha}\ }
 \alpha_{12},
 \qquad
 E_{21/22}^{(9)}:I_\alpha\longrightarrow\mathbb R,
 \quad E_{21/22}^{(9)}(\xi_9)=0,
 \quad (E_{21/22}^{(9)})'(\xi_9)\ne0.
 \label{eq:l9s102:alpha-dos-mapas}
\end{equation}
La clausura entera publica primero la ventana finita \(\alpha_{12}\) y su
familia compatible de prolongaciones; su límite inverso es
\(\alpha_{\mathrm{HMT}}\). El segundo lector publica \(\xi_9\) como testigo
finito de orden nueve y, por iteración del mismo transporte graduado del TPK,
construye el sistema compatible completo
\(E_{21/22}^{(6+3m)}\). Los dos sistemas inversos determinan en cada
profundidad el mismo cilindro superviviente \(C_N^\alpha\); sus diámetros
tienden a cero. Por tanto, la intersección es unitaria y el teorema de las
dos vías establece
\begin{equation}
 \alpha_A:=\varprojlim_N\rho_N(\alpha_{12}),
 \qquad
 E_{21/22}:=\varprojlim_m E_{21/22}^{(6+3m)},
 \qquad
 \alpha_B:=\operatorname*{arg\,zero}_{x\in I_\alpha}E_{21/22}(x),
 \qquad
 \boxed{\alpha_A=\alpha_B=\alpha_{\mathrm{HMT}}}.
 \label{eq:l9s102:alpha-identidad-limite-dos-vias}
\end{equation}
Ninguna vía introduce un valor objetivo como dato de entrada; la
identificación metrológica pertenece únicamente al reconocimiento
convencional posterior. La igualdad
\(\operatorname{Tr}_9(\xi_9^{-1})=
\operatorname{Tr}_9(\alpha_{12}^{-1})=137.035999084\) es un certificado
finito de la identidad anterior, nunca su definición ni su alcance máximo.
La ley de prolongación no se añade como dato: es la iteración del transporte
graduado ya producido por APP--TRIT--TPK, con retorno de fase, avance de
memoria y truncamientos compatibles.

\subsection{Núcleo analítico \texorpdfstring{\(P_6\)}{P6} y prolongación exacta de orden \texorpdfstring{\(9\)}{9}}
\label{sec:l9s102:alpha-p6}

La caracterización de vacancias emplea
\begin{equation}
 \begin{aligned}
 P_6(x)={}&2\pi_{\mathrm{HMT}}x-\frac74x^2
 +\frac{x^3}{2\pi_{\mathrm{HMT}}}
 +\frac{x^4}{20}-\frac{2x^5}{21}-\frac{x^6}{46},\\
 P_6(x)={}&\log_{10}\frac{10}{9}.
 \end{aligned}
 \label{eq:l9s102:p6}
\end{equation}
La raíz positiva certificada de este jet satisface
\begin{equation}
 x_6^{-1}=137.03599908400002702009\ldots .
 \label{eq:l9s102:raíz-p6}
\end{equation}
El polinomio \(P_6\) proporciona el jet inicial de la caracterización
analítica. La holonomía nonádica prolonga exactamente sus grados \(7,8,9\) y
la misma regla graduada se itera sobre los grados sucesivos. El cociente de
jets es el primer certificado finito de esa prolongación coinductiva; no
selecciona ni sustituye su sección límite.

\subsection{Genealogía APP--TRIT--TPK de los invariantes \texorpdfstring{\(120,45,11\)}{120, 45, 11} y \texorpdfstring{\(495\)}{495}}
\label{sec:l9s102:alpha-genealogia-7-9}

La semicorona de fase del TPK construye
\begin{equation}
 |\mathcal F_8^{\mathrm{ph}}|=120,
 \qquad |L_{\mathrm{ph}}|=45.
 \label{eq:l9s102:alpha-120-45}
\end{equation}
APP aporta el bloque central
\begin{equation}
 B_c=\begin{pmatrix}7&2\\2&7\end{pmatrix},
 \qquad
 \det B_c=45,
 \qquad
 \operatorname{Spec}(B_c)=\{9,5\}.
 \label{eq:l9s102:alpha-bc}
\end{equation}
El extractor transversal dodecafásico proporciona
\begin{equation}
 r_{\mathrm{dod}}=\operatorname{rank}(D_3,D_4)=11,
 \qquad
 (D_3K)_1=-495,
 \qquad
 495=45\cdot11=\det(B_c)\,r_{\mathrm{dod}}.
 \label{eq:l9s102:alpha-495}
\end{equation}
Los coeficientes del jet prolongado proceden, por tanto, de tres
invariantes previos de APP--TRIT--TPK:
\begin{equation}
 -\frac1{120}\longmapsto\frac1{45}
 \longmapsto\frac2{45\cdot11}.
 \label{eq:l9s102:alpha-coeficientes}
\end{equation}

\subsection{Resolvente nilpotente \texorpdfstring{\((I-N)^{-1}\)}{(I-N)^-1} y construcción exacta de la cola \texorpdfstring{\(T_{7:9}\)}{T7:9}}
\label{sec:l9s102:alpha-resolvente}

Sobre la base \(e_7,e_8,e_9\), definimos
\begin{equation}
 Ne_7=-\frac83e_8,
 \qquad
 Ne_8=\frac2{11}e_9,
 \qquad
 Ne_9=0,
 \qquad
 v_7=-\frac1{120}e_7.
 \label{eq:l9s102:alpha-nilpotente}
\end{equation}
Como \(N^3=0\), su resolvente finito actua exactamente por
\begin{equation}
 (I-N)^{-1}v_7
 =v_7+Nv_7+N^2v_7
 =-\frac{e_7}{120}+\frac{e_8}{45}+\frac{2e_9}{495}.
 \label{eq:l9s102:alpha-resolvente-finito}
\end{equation}
La evaluación \(e_n\mapsto x^n\) produce
\begin{equation}
 \boxed{
 T_{7:9}(x)=-\frac{x^7}{120}+\frac{x^8}{45}
 +\frac{2x^9}{495}.}
 \label{eq:l9s102:alpha-t79}
\end{equation}

\subsection{Jet nonádico \texorpdfstring{\(P_9\)}{P9}, raíz prolongada y escala de residuos certificados}
\label{sec:l9s102:alpha-jet-p9}

El jet nonádico de orden \(9\) queda definido por
\begin{equation}
 P_9(x)=P_6(x)+T_{7:9}(x),
 \qquad
 P_9(x)=\log_{10}\frac{10}{9}.
 \label{eq:l9s102:alpha-p9}
\end{equation}
\begin{theorem}[Existencia, simplicidad y unicidad de la segunda salida analítica]
\label{thm:l9s102:raiz-p9-unica}
El operador \(E_{21/22}^{(9)}\) posee una única raíz positiva \(\xi_9\) en
\(I_\alpha=(0,10^{-2})\). Está encerrada por el intervalo racional exacto
\begin{equation}
 \frac{729735256928380099}{10^{20}}
 <\xi_9<
 \frac{729735256928380100}{10^{20}},
 \label{eq:l9s102:encierro-raiz-p9}
\end{equation}
y su inversa satisface
\begin{equation}
 \xi_9^{-1}
 =137.0359990840000000000000111926291900\ldots .
 \label{eq:l9s102:alpha-raíz-p9}
\end{equation}
\end{theorem}

\begin{proof}
Todos los coeficientes de \(P_9\) se extraen de los invariantes anteriores a
la resolución. Para \(0\leq x\leq10^{-2}\), usando
\(3.14<\pi_{\rm HMT}<3.15\), se obtiene
\[
 (E_{21/22}^{(9)})'(x)
 \geq 2\pi_{\rm HMT}-\frac72\,10^{-2}
 -\frac{10}{21}\,10^{-8}-\frac6{46}\,10^{-10}
 -\frac7{120}\,10^{-12}>6.24.
\]
Las evaluaciones dirigidas en los extremos de
\cref{eq:l9s102:encierro-raiz-p9} tienen signos opuestos. El teorema del
valor intermedio proporciona la existencia; la cota derivativa proporciona
la unicidad y la simplicidad. La división intervalar produce la expansión de
\(\xi_9^{-1}\), que constituye el testigo de orden nueve del sistema
compatible completo.
\end{proof}

\begin{table}[htbp]
\centering
\caption{Descenso certificado del residuo de la caracterización analítica}
\label{tab:l9s102:alpha-residuos}
\begin{tabular}{@{}lc@{}}
\toprule
Núcleo evaluado en \(\alpha_{12}\) & Residuo \\
\midrule
\(P_6\) & \(9.0038886\times10^{-18}\) \\
\(P_6-x^7/120\) & \(-1.7893115\times10^{-19}\) \\
\(P_6-x^7/120+x^8/45\) & \(-2.3708601\times10^{-22}\) \\
\(P_9\) & \(3.7297132\times10^{-27}\) \\
\bottomrule
\end{tabular}
\end{table}

\subsection{Operador dodecafásico alternativo \texorpdfstring{\(A_{\mathrm{dod}}\)}{A_dod} y control de unicidad del lector tipado}
\label{sec:l9s102:alpha-control}

El operador de control
\begin{equation}
 A_{\mathrm{dod}}
 =I-\frac14(D_3^*D_3+D_4^*D_4)
 \label{eq:l9s102:alpha-control-alternativo}
\end{equation}
produce una sucesión diferente de coeficientes. Este control demuestra que
la simetría dodecafásica aislada resulta insuficiente para seleccionar la
prolongación \(T_{7:9}\); la selección pertenece a la composición tipada del
lector \(\Gamma_{E21}^{(9)}\) con el estado enriquecido. La diferencia entre
ambos operadores certifica la necesidad de conservar dominio, orientación,
graduación y memoria, y prueba la especificidad del lector finito. La cola
analítica pertenece a la iteración del lector tipado completo, no a la
simetría aislada de este operador de control.

\subsection{Alcance exacto del operador nilpotente y exclusión de colas ajenas al transporte TPK}
\label{sec:l9s102:alpha-alcance}

La aritmética de \(120,45,11,495\), el operador nilpotente, el resolvente y
la identidad \cref{eq:l9s102:alpha-t79} poseen exactitud interna. La raíz
positiva simple y única de \(E_{21/22}^{(9)}\) constituye el testigo de orden
nueve de la segunda lectura interna. La concordancia de \(P_9\) con la
ventana dodecafásica certifica el primer cuadrado finito del sistema
compatible. La unicidad canónica de la prolongación \(7{:}9\) queda tipada
por el morfismo
\begin{equation}
 \Gamma_{E21}^{(9)}:
 X_{\mathrm{APP\text{-}TRIT\text{-}TPK}}
 \longrightarrow \mathbb Q(\pi_{\mathrm{HMT}})[x]/(x^{10}),
 \label{eq:l9s102:alpha-gamma-e21}
\end{equation}
cuya acción determina conjuntamente la semilla orientada, las dos
transferencias, la graduación \(7,8,9\), el cociente de jets
\(F^7/F^{10}\) y su transporte bajo el retorno de fase.

Sea \(K=\mathbb Q(\pi_{\rm HMT},\log_{10}(10/9))\) y sea
\(j_9:K[[x]]\to K[x]/(x^{10})\) el truncamiento. Entonces
\begin{equation}
 \ker j_9=x^{10}K[[x]],
 \qquad
 E_h(x):=E_{21/22}^{(9)}(x)+x^{10}h(x),
 \qquad h\in K[[x]].
\label{eq:l9s102:alpha-limite-jets-e21}
\end{equation}
Todos los \(E_h\) poseen el mismo jet certificado. Sin embargo,
\(E_0\) y \(E_{1/729}\) tienen raíces distintas en \(I_\alpha\), pues
\(E_{1/729}(\xi_9)=\xi_9^{10}/729>0\). Esta observación demuestra únicamente
que un jet aislado no determina una cola formal. No afecta a la segunda vía
HMT, porque ésta no permite elegir libremente \(h\): la recurrencia del
transporte TPK construye la única familia compatible
\(E_{21/22}^{(6+3m)}\), y sus truncamientos conmutan con los de la vía
dodecafásica. Elegir \(h\) mediante el valor objetivo sería circular y queda
excluido; elegirlo mediante el transporte generado es precisamente la
operación ya demostrada.

El falsador separa así tres afirmaciones: invalida el certificado finito si
fallan la monotonía, el cambio de signo o el aislamiento de \(\xi_9\);
invalida la prolongación \(7{:}9\) si sus coeficientes dejan de factorizar por
\(\Gamma_{E21}^{(9)}\); e invalida la identidad completa si falla la
compatibilidad de algún cuadrado del sistema inverso o si los diámetros de
los cilindros supervivientes no tienden a cero. Ninguna de estas condiciones
introduce un morfismo pendiente: son los falsadores materiales de la
construcción que establece
\(\alpha_A=\alpha_B=\alpha_{\mathrm{HMT}}\).
